\section{Gaussian process background model and statistical method}
\label{sec:gpr-stats}

For each data sample and mass hypothesis, the analysis removes a narrow region around
the tested mass from the background training set. The nominal observed-limit workflow
uses a $2.25\,\sigma_m$ blind and extraction half-width, where $\sigma_m$ is the
data-set-specific mass resolution. The Gaussian process is trained on the remaining
sideband bins across the full histogram range. Its posterior prediction inside the
blind window supplies both the background mean and the bin-to-bin covariance used in
the local likelihood.

\begin{figure*}[t]
\centering
\begin{minipage}{0.47\textwidth}
\centering
\ifigure{\linewidth}{gpr_intro_cartoon.png}\\
(a) Sideband-trained GPR interpolation.
\end{minipage}
\hfill
\begin{minipage}{0.47\textwidth}
\centering
\ifigure{\linewidth}{gpr_training_steps.png}\\
(b) Log-count training and blind-window prediction.
\end{minipage}
\caption{GPR construction used in the prompt-resonance search. For each mass
hypothesis, the bins near the signal are excluded from training. The trained Gaussian
process predicts a background mean and covariance in the excluded window, which are
then passed to the profile likelihood.}
\label{fig:gpr-method}
\end{figure*}

The fit is performed in transformed coordinates,
\begin{equation}
x=\log m,\qquad z=\log y,
\end{equation}
for occupied mass bins with count $y$. This transformation is useful because the
prompt spectrum falls steeply and spans a wide dynamic range. The baseline kernel is a
single constant-times-radial-basis-function kernel,
\begin{equation}
k(x,x')=C\exp\left[-\frac{(x-x')^2}{2\ell^2}\right],
\end{equation}
with the length-scale bounds tied to the detector resolution in log-mass space. This
choice follows the smooth-background GPR strategy used in high-energy physics
applications while keeping the HPS baseline model intentionally simple
\cite{RasmussenWilliams2006,FrateGP2017,Gandrakota2023,BarrLiu2025}. Alternative
kernels are treated as robustness studies rather than as competing mainline methods.

After conditioning on the sidebands, the Gaussian process posterior in log-count space
is mapped back to count space using the standard lognormal moment relations. The
result is a background mean vector $\bm{b}_d(m)$ and covariance matrix
$\bm{C}_d(m)$ for the blinded bins of data set $d$ at mass $m$. The covariance is
factorized as
\begin{equation}
\bm{C}_d=L_dL_d^T,
\end{equation}
and propagated as correlated Gaussian-constrained nuisance parameters
$L_d\bm{\theta}_d$ in the likelihood. This is the standard high-energy physics way to
encode correlated background uncertainty in a profile likelihood
\cite{HistFactory,Cowan2011}.

The signal template is a detector-resolution Gaussian normalized over the full
histogram range and then restricted to the blinded bins,
\begin{equation}
w_{d,i}(m)=
\int_{e_i}^{e_{i+1}}
\frac{1}{\sqrt{2\pi}\sigma_d(m)}
\exp\left[-\frac{(m'-m)^2}{2\sigma_d^2(m)}\right]dm',
\end{equation}
for bins $i$ in the local blind window. The template is not renormalized after this
restriction. The fitted amplitude $A_d$ is therefore the total local Gaussian signal
yield, while only the fraction of that Gaussian inside the blind window contributes to
the local likelihood. This convention preserves the physical meaning of the signal
yield and avoids silently discarding signal tails.

For a fixed mass, define the expected bin count
\begin{equation}
\lambda_{d,i}(A_d,\bm{\theta}_d)=
b_{d,i}+(L_d\bm{\theta}_d)_i+A_dw_{d,i}.
\end{equation}
The single-sample likelihood is then
\begin{equation}
\mathcal{L}_d(A_d,\bm{\theta}_d)=
\left[
\prod_i \mathrm{Pois}\left(n_{d,i}\middle|\lambda_{d,i}\right)
\right]
\exp\left(-\frac{1}{2}\bm{\theta}_d^T\bm{\theta}_d\right).
\label{eq:single-likelihood}
\end{equation}
The local discovery test uses the one-sided profile-likelihood statistic for
$A_d=0$, with the Gaussian-equivalent local significance obtained using the asymptotic
mapping of Cowan \textit{et al.} \cite{Cowan2011}. Upper limits use the bounded
one-sided statistic $\tilde q_A$ and the modified-frequentist construction
\begin{equation}
\CLs(A;m)=\frac{\mathrm{CL}_{s+b}(A;m)}{\mathrm{CL}_{b}(A;m)}.
\end{equation}
The quoted upper limit is the value for which $\CLs=0.05$ for a 95\% confidence-level
limit, or $\CLs=0.10$ for a 90\% confidence-level limit. Dividing by
$\mathrm{CL}_b$ protects against quoting overly strong exclusions from downward
background fluctuations \cite{Junk1999,Read2002}.

Observed and expected limits differ by the data vector or ensemble used in this
construction. Observed limits use the actual blinded-bin counts. Expected bands use
background-only pseudoexperiments, or the equivalent asymptotic background-only
reference, conditioned on the sideband-trained GPR prediction. The local density
$\rho_d(m)$ used in Eq.~\eqref{eq:kd} is a normalization input for converting yield to
$\eps^2$; it is not the definition of observed versus expected limits.

When multiple data samples are active at the same mass, the combined likelihood is the
product of the sample likelihoods,
\begin{equation}
\mathcal{L}_{\mathrm{comb}}(\eps^2,\{\bm{\theta}_d\})
=
\prod_{d\in\mathcal{D}(m)}
\mathcal{L}_d\left(A_d=K_d^{\mathrm{eff}}(m)\eps^2,\bm{\theta}_d\right),
\label{eq:combined-likelihood}
\end{equation}
where $\mathcal{D}(m)$ is the active set of samples. The parameter of interest is the
shared coupling $\eps^2$. The nuisance parameters, background predictions, mass
resolutions, and radiative normalizations remain data-set-specific. Thus the combined
limit is a direct shared-coupling inference, not an average, product, or envelope of
individual upper-limit curves.

The local look-elsewhere effect is handled separately from the fixed-mass test. The
preferred global discovery calibration is the background-only distribution of the
maximum discovery statistic over the mass scan,
\begin{equation}
p_{\mathrm{global}}=
\frac{1}{N_{\mathrm{toys}}}
\sum_{t=1}^{N_{\mathrm{toys}}}
\mathbf{1}\left[q_{0,\max}^{(t)}\ge q_{0,\max}^{\mathrm{obs}}\right].
\end{equation}
Until that scan-level toy distribution is frozen for the final production result,
effective-trials overlays are used only as review diagnostics.
